\documentclass[11pt]{article}
\usepackage[a4paper,margin=1in]{geometry}
\usepackage{fontspec}
\usepackage[boldfont=false]{xeCJK}
\setCJKmainfont{FandolSong-Regular.otf}
\usepackage{amsmath}
\usepackage{amssymb}
\usepackage{array}
\usepackage{booktabs}
\usepackage{graphicx}
\usepackage{adjustbox}
\usepackage{enumitem}
\setlist[itemize]{topsep=2pt,partopsep=0pt,parsep=0pt,itemsep=2pt,leftmargin=1.4em}
\setlist[enumerate]{topsep=2pt,partopsep=0pt,parsep=0pt,itemsep=2pt,leftmargin=1.6em}
\usepackage{parskip}
\usepackage{fvextra}
\fvset{breaklines=true,breakanywhere=true,fontsize=\small}
\setlength{\parindent}{0pt}

\begin{document}

\begin{center}
  {\LARGE\bfseries Experimental techniques and chemical analysis}\\[0.35em]
  {\large Igcse Chemistry}
\end{center}
\vspace{0.6em}

\section*{Experimental design}

\subsection*{Apparatus}

You should know which piece of apparatus to use for each measurement:

\begin{center}
\adjustbox{max width=\linewidth}{%
\begin{tabular}{ll}
\toprule
\textbf{Measurement} & \textbf{Apparatus} \\
\midrule
time & \textbf{stop-watch} 秒表 \\
temperature & \textbf{thermometer} 温度计 \\
mass & \textbf{balance} 天平 \\
volume (accurate) & \textbf{burette} 滴定管 or \textbf{volumetric pipette} 移液管 \\
volume (rough) & \textbf{measuring cylinder} 量筒 \\
volume of a gas & \textbf{gas syringe} 注射器 \\
\bottomrule
\end{tabular}}
\end{center}

\subsection*{Key words}

These words are used throughout practical chemistry:

\begin{itemize}
\item A \textbf{solvent} 溶剂 is a substance that dissolves a solute.

\item A \textbf{solute} 溶质 is the substance that dissolves in the solvent.

\item A \textbf{solution} 溶液 is a mixture of one or more solutes dissolved in a solvent.

\item A \textbf{saturated solution} 饱和溶液 holds the maximum amount of solute that will dissolve at a given temperature.

\item A \textbf{residue} 残渣 is the substance left behind after evaporation, distillation or filtration.

\item A \textbf{filtrate} 滤液 is the liquid that has passed through a filter.

\end{itemize}

\section*{Acid–base titrations}

A \textbf{titration} 滴定 finds the exact volume of one solution that reacts with another.

\begin{itemize}
\item Use a \textbf{volumetric pipette} to measure a fixed volume of one solution into a flask.

\item Add a few drops of a suitable \textbf{indicator} 指示剂.

\item Add the other solution from a \textbf{burette}, slowly, until the colour just changes.

\end{itemize}

The \textbf{end-point} 终点 is the moment the indicator changes colour, which shows the reaction is exactly complete. You read the burette to find the volume added.

\section*{Chromatography}

\textbf{Paper chromatography} 纸色谱法 separates a mixture of soluble substances. You put a spot of the mixture near the bottom of the paper, then stand the paper in a solvent. As the solvent rises up the paper, the substances move different distances, so they separate.

\begin{itemize}
\item For \textbf{coloured} substances you can see the spots directly.

\item For \textbf{colourless} substances you must spray a \textbf{locating agent} 显色剂 to make the spots show up.

\end{itemize}

The finished paper is a \textbf{chromatogram} 色谱图. You can use it to:

\begin{itemize}
\item identify an unknown substance by comparing it with known substances;

\item tell if a substance is pure (a pure substance gives only one spot).

\end{itemize}

You can also calculate the $R_{\text{f}}$ value of a spot:

\[
R_{\text{f}} = \frac{\text{distance moved by the substance}}{\text{distance moved by the solvent}}
\]

\section*{Separation and purification}

You choose a method based on the mixture:

\begin{center}
\adjustbox{max width=\linewidth}{%
\begin{tabular}{ll}
\toprule
\textbf{Method} & \textbf{Used to separate} \\
\midrule
dissolving in a suitable solvent, then \textbf{filtration} 过滤 & an insoluble solid from a liquid \\
\textbf{crystallisation} 结晶 & a soluble solid from its solution \\
\textbf{simple distillation} 蒸馏 & a solvent (the liquid) from a solution \\
\textbf{fractional distillation} 分馏 & two or more liquids with different boiling points \\
\bottomrule
\end{tabular}}
\end{center}

You can check the purity of a substance using its \textbf{melting point} 熔点 and \textbf{boiling point} 沸点: a pure substance melts and boils at sharp, fixed temperatures, while impurities lower the melting point and raise the boiling point.

\section*{Identifying ions and gases}

\subsection*{Tests for anions}

An \textbf{anion} 阴离子 is a negative ion.

\begin{center}
\adjustbox{max width=\linewidth}{%
\begin{tabular}{lll}
\toprule
\textbf{Anion} & \textbf{Test} & \textbf{Result} \\
\midrule
\textbf{carbonate} 碳酸盐 ($\text{CO}_3^{2-}$) & add dilute acid & fizzes; the gas turns limewater milky (\textbf{carbon dioxide} 二氧化碳) \\
\textbf{chloride} 氯化物 ($\text{Cl}^-$) & add dilute \textbf{nitric acid} 硝酸, then \textbf{silver nitrate} 硝酸银 & white precipitate \\
\textbf{bromide} 溴化物 ($\text{Br}^-$) & add dilute nitric acid, then silver nitrate & cream precipitate \\
\textbf{iodide} 碘化物 ($\text{I}^-$) & add dilute nitric acid, then silver nitrate & yellow precipitate \\
\textbf{nitrate} 硝酸盐 ($\text{NO}_3^-$) & add \textbf{aluminium} 铝 foil and \textbf{sodium hydroxide} 氢氧化钠, warm & \textbf{ammonia} 氨气 gas given off \\
\textbf{sulfate} 硫酸盐 ($\text{SO}_4^{2-}$) & add dilute nitric acid, then \textbf{barium nitrate} 硝酸钡 & white precipitate \\
\textbf{sulfite} 亚硫酸盐 ($\text{SO}_3^{2-}$) & add acidified \textbf{potassium manganate(VII)} 高锰酸钾 & purple colour fades \\
\bottomrule
\end{tabular}}
\end{center}

\subsection*{Tests for cations}

A \textbf{cation} 阳离子 is a positive ion. Add aqueous sodium hydroxide, or aqueous ammonia, and look at the \textbf{precipitate} 沉淀 formed.

\begin{center}
\adjustbox{max width=\linewidth}{%
\begin{tabular}{lll}
\toprule
\textbf{Cation} & \textbf{With sodium hydroxide} & \textbf{With aqueous ammonia} \\
\midrule
aluminium ($\text{Al}^{3+}$) & white, dissolves in excess & white, stays \\
\textbf{ammonium} 铵 ($\text{NH}_4^+$) & ammonia gas when warmed & — \\
\textbf{calcium} 钙 ($\text{Ca}^{2+}$) & white, stays & no precipitate \\
\textbf{chromium}(III) 铬 ($\text{Cr}^{3+}$) & green, dissolves in excess & green, stays \\
\textbf{copper}(II) 铜 ($\text{Cu}^{2+}$) & light blue, stays & light blue, dissolves to deep blue \\
\textbf{iron}(II) 铁 ($\text{Fe}^{2+}$) & green, stays & green, stays \\
iron(III) ($\text{Fe}^{3+}$) & red-brown, stays & red-brown, stays \\
\textbf{zinc} 锌 ($\text{Zn}^{2+}$) & white, dissolves in excess & white, dissolves in excess \\
\bottomrule
\end{tabular}}
\end{center}

\subsection*{Tests for gases}

\begin{center}
\adjustbox{max width=\linewidth}{%
\begin{tabular}{lll}
\toprule
\textbf{Gas} & \textbf{Test} & \textbf{Result} \\
\midrule
ammonia ($\text{NH}_3$) & damp red \textbf{litmus} 石蕊 paper & turns blue \\
carbon dioxide ($\text{CO}_2$) & bubble through \textbf{limewater} 石灰水 & turns milky \\
\textbf{chlorine} 氯气 ($\text{Cl}_2$) & damp litmus paper & bleached white \\
\textbf{hydrogen} 氢气 ($\text{H}_2$) & a lighted splint & burns with a squeaky pop \\
\textbf{oxygen} 氧气 ($\text{O}_2$) & a glowing splint & relights \\
sulfur dioxide ($\text{SO}_2$) & acidified potassium manganate(VII) & purple colour fades \\
\bottomrule
\end{tabular}}
\end{center}

\subsection*{Flame tests}

A \textbf{flame test} 焰色试验 identifies some metal cations by the colour they give to a flame:

\begin{center}
\adjustbox{max width=\linewidth}{%
\begin{tabular}{ll}
\toprule
\textbf{Cation} & \textbf{Flame colour} \\
\midrule
\textbf{lithium} 锂 ($\text{Li}^+$) & red \\
\textbf{sodium} 钠 ($\text{Na}^+$) & yellow \\
\textbf{potassium} 钾 ($\text{K}^+$) & lilac (purple) \\
calcium ($\text{Ca}^{2+}$) & orange-red \\
\textbf{barium} 钡 ($\text{Ba}^{2+}$) & light green \\
copper(II) ($\text{Cu}^{2+}$) & blue-green \\
\bottomrule
\end{tabular}}
\end{center}


\end{document}
